\documentclass[SE,lsstdraft,authoryear,toc]{lsstdoc}
\input{meta}

% Package imports go here.

% Local commands go here.

%If you want glossaries
%\input{aglossary.tex}
%\makeglossaries

\title{Thermal Predictors of Open-Loop Focus Drift at the Vera C. Rubin Observatory}

% This can write metadata into the PDF.
% Update keywords and author information as necessary.
\hypersetup{
    pdftitle={Thermal Predictors of Open-Loop Focus Drift at the Vera C. Rubin Observatory},
    pdfauthor={ChristineCorry},
    pdfkeywords={}
}

% Optional subtitle
% \setDocSubtitle{A subtitle}

\input{authors}

\setDocRef{SOTN-005}
\setDocUpstreamLocation{\url{https://github.com/lsst-so/sotn-005}}

\date{\vcsDate}

% Optional: name of the document's curator
% \setDocCurator{The Curator of this Document}

\setDocAbstract{%
The Vera C. Rubin Observatory’s Active Optics System corrects slowly varying, thermally induced defocus through closed-loop active control and, ultimately, an open-loop look-up table calibrated against bulk temperature. However, rapid residual focus drifts, quantified by the Zernike coefficient $Z_4$, are observed after LUT correction on timescales of $\sim 30$ minutes. These drifts are suspected to arise from rapidly evolving thermo-optical coupling, motivating a systematic characterization of residual thermal effects on focus stability. We analyze 23 open-loop stability image sequences from November and December 2025, each consisting of 40 exposures acquired over $\sim$28 minutes with the hexapod held fixed. Temperature data from 13 sensors spanning the enclosure air, telescope structure, mirror glass, and hexapod assemblies are queried and mean-centered to isolate thermal dynamics from bulk nightly cooling. We compute Pearson correlations between $Z_4$ and individual sensors, physically motivated pairwise temperature gradient proxies, and PCA-derived thermal modes, repeating all analyses at the slope timescale. Individual sensor and temperature gradient proxy correlations are weak (max $r \simeq 0.3$) and no sensor maintains a stable multivariate coefficient across runs. Top sensor thermal correlations are stronger in runs classified as smooth, consistent with a thermal signal that is present but frequently obscured by measurement noise and other dynamical effects. Correlations among the $Z_4$ measurements themselves further complicate the interpretation, indicating that apparent sensor relationships may partly reflect shared temporal structure rather than direct causal coupling. These results suggest that thermal-optical coupling at the Rubin Observatory is distributed across multiple mechanisms and timescales and cannot be captured by simple linear relationships with individual temperature sensors. Future work will require additional temperature sensing, a larger sample of open-loop runs, and nonlinear or multivariate dynamical models to separate true thermal drivers from correlated focus variability.
}

% Change history defined here.
% Order: oldest first.
% Fields: VERSION, DATE, DESCRIPTION, OWNER NAME.
% See LPM-51 for version number policy.
\setDocChangeRecord{%
  \addtohist{1}{YYYY-MM-DD}{Unreleased.}{Corry}
}


\begin{document}

% Create the title page.
\maketitle
% Frequently for a technote we do not want a title page  uncomment this to remove the title page and changelog.
% use \mkshorttitle to remove the extra pages

% ADD CONTENT HERE
% You can also use the \input command to include several content files.

\appendix

\section{Acknowledgements}

This material is based upon work supported in part by the National Science Foundation through Cooperative Agreements AST-1258333 and AST-2241526 and Cooperative Support Agreements AST-1202910 and AST-2211468 managed by the Association of Universities for Research in Astronomy (AURA), and the Department of Energy under Contract No.\ DE-AC02-76SF00515 with the SLAC National Accelerator Laboratory managed by Stanford University.
Additional Rubin Observatory funding comes from private donations, grants to universities, and in-kind support from LSST-DA Institutional Members.

% Include all the relevant bib files.
% https://lsst-texmf.lsst.io/lsstdoc.html#bibliographies
\section{References} \label{sec:bib}
\renewcommand{\refname}{} % Suppress default Bibliography section
\bibliography{local,lsst,lsst-dm,refs_ads,refs,books}

% Make sure lsst-texmf/bin/generateAcronyms.py is in your path
\section{Acronyms} \label{sec:acronyms}
\input{acronyms.tex}
% If you want glossary uncomment below -- comment out the two lines above
%\printglossaries





\end{document}
